%%%%%%%%%%%%%%%%%%%%%%%%%%%%%%%%%%%%%%%%%%%%%%%%%%%%%%%%%%%%%%%%%%%%%%%%%%%%%%%
%%  TKS_Cognitive_Warfare_Stack.tex
%%  The Integrated Framework: Individual -> Social -> Cultural -> Civilizational
%%
%%  Part of TKS v7.4 - December 2025
%%  ADDITIVE to v7.4 canon
%%
%%  This document provides comprehensive technical documentation for the
%%  multi-level prediction and influence framework that emerges from TKS.
%%  Presented without advocacy - documenting what exists, not what should be.
%%%%%%%%%%%%%%%%%%%%%%%%%%%%%%%%%%%%%%%%%%%%%%%%%%%%%%%%%%%%%%%%%%%%%%%%%%%%%%%

\documentclass[11pt,twoside]{book}

%% ============================================================================
%% PACKAGES
%% ============================================================================
\usepackage[utf8]{inputenc}
\usepackage[T1]{fontenc}
\usepackage[margin=1in,inner=1.2in,outer=1in]{geometry}
\usepackage{amsmath,amssymb,amsthm}
\usepackage{mathtools}
\usepackage{stmaryrd}
\usepackage{enumitem}
\usepackage{booktabs}
\usepackage{longtable}
\usepackage{tabularx}
\usepackage{array}
\usepackage{tcolorbox}
\tcbuselibrary{breakable,skins}
\usepackage{tikz-cd}
\usepackage{tikz}
\usetikzlibrary{positioning,shapes,arrows.meta,calc}
\usepackage{xcolor}
\usepackage{hyperref}
\usepackage{ragged2e}

%% ============================================================================
%% FORMATTING SETTINGS
%% ============================================================================
\setlength{\parindent}{0pt}
\setlength{\parskip}{0.5em}
\renewcommand{\arraystretch}{1.2}

%% ============================================================================
%% CUSTOM COMMANDS (from TKS canonical)
%% ============================================================================
\newcommand{\noe}[1]{\nu_{#1}}
\newcommand{\Ideas}{\mathcal{I}}
\newcommand{\Worlds}{\mathcal{W}}
\newcommand{\Elements}{\mathcal{E}}
\newcommand{\Noetics}{\mathcal{N}}
\newcommand{\Foundations}{\mathcal{F}}
\newcommand{\Acquisitions}{\mathcal{A}}
\newcommand{\Frac}{\Phi}
\newcommand{\TransFrac}[1]{\Phi^{#1}}
\newcommand{\MultiFrac}[2]{\Phi^{#1}_{#2}}
\newcommand{\RPM}{\mathsf{RPM}}
\newcommand{\ACBE}{\mathsf{ACBE}}
\newcommand{\HausdorffDim}{\dim_H}
\newcommand{\FractalSpectrum}{\sigma_\Phi}
\newcommand{\fix}{\mathrm{fix}}
\newcommand{\sem}[1]{\llbracket #1 \rrbracket}
\newcommand{\Tadd}{\oplus_T}
\newcommand{\Tsub}{\ominus_T}
\newcommand{\Tmul}{\otimes_T}
\newcommand{\Tdiv}{\oslash_T}
\newcommand{\Wadd}{\oplus_W}
\newcommand{\Wsub}{\ominus_W}
\newcommand{\Failed}{\mathtt{Failed}}

%% New commands for Cognitive Warfare Stack
\newcommand{\Level}[1]{\mathfrak{L}_{#1}}
\newcommand{\StackFunctor}{\mathcal{S}}
\newcommand{\Individual}{\mathfrak{I}}
\newcommand{\Social}{\mathfrak{S}}
\newcommand{\Cultural}{\mathfrak{C}}
\newcommand{\Civilizational}{\mathfrak{Z}}
\newcommand{\PsychAttractor}{\mathcal{A}_\psi}
\newcommand{\GroupFrac}{\Phi_{\mathrm{group}}}
\newcommand{\Overton}{\mathcal{O}}
\newcommand{\Memetic}{\mathcal{M}}
\newcommand{\Hyperstition}{\mathcal{H}}
\newcommand{\Trajectory}{\mathcal{T}}
\newcommand{\Vulnerability}{\mathcal{V}}
\newcommand{\Cascade}{\mathcal{C}}
\newcommand{\Aggregate}[2]{\mathbf{Agg}_{#1}^{#2}}
\newcommand{\Project}[2]{\mathbf{Proj}_{#1}^{#2}}

%% Theorems
\newtheorem{definition}{Definition}[chapter]
\newtheorem{theorem}{Theorem}[chapter]
\newtheorem{proposition}{Proposition}[chapter]
\newtheorem{lemma}{Lemma}[chapter]
\newtheorem{corollary}{Corollary}[chapter]
\newtheorem{example}{Example}[chapter]
\newtheorem{remark}{Remark}[chapter]
\newtheorem{axiom}{Axiom}[chapter]

%% ============================================================================
%% DOCUMENT
%% ============================================================================
\begin{document}

\title{\textbf{TKS Cognitive Warfare Stack} \\[0.5em]
\Large The Integrated Multi-Level Framework \\
for Prediction, Modeling, and Influence \\[0.3em]
\normalsize Individual $\to$ Social $\to$ Cultural $\to$ Civilizational}
\author{TKS v7.4 Formalization Project}
\date{December 2025}
\maketitle

\tableofcontents

%% ============================================================================
%% PREFACE: DOCUMENTATION PURPOSE
%% ============================================================================
\chapter*{Preface: On Documentation}
\addcontentsline{toc}{chapter}{Preface: On Documentation}

\begin{tcolorbox}[colback=gray!10!white,colframe=gray!75!black,title=\textbf{Scope and Intent},breakable]
This document provides comprehensive technical documentation for what we term the
\textbf{TKS Cognitive Warfare Stack}---the integrated framework showing how TKS
prediction and influence capabilities layer together across scales, from individual
psychology to civilization-level dynamics.

\textbf{This document is descriptive, not prescriptive.} The mathematics and
mechanisms described herein exist regardless of whether they are documented.
Advertising firms, political campaigns, intelligence agencies, therapeutic
practitioners, and social engineers have employed analogous methods throughout
history. TKS provides formal mathematical language for what has always been done
informally.

The question is not whether such capabilities exist---they do. The question is
who understands them, who uses them, and for what purposes. Opacity benefits
those with existing power and knowledge; transparency enables informed consent
and counter-measures.
\end{tcolorbox}

%% ============================================================================
%% PART I: MATHEMATICAL ARCHITECTURE
%% ============================================================================
\part{Mathematical Architecture of the Stack}

\chapter{The Four-Level Structure}
\label{ch:four-level-structure}

\section{Overview: Levels and Their Domains}

The TKS Cognitive Warfare Stack comprises four hierarchical levels, each operating
at a distinct scale of consciousness organization:

\begin{center}
\begin{tikzpicture}[
    level/.style={rectangle, draw, rounded corners, minimum width=10cm, minimum height=1.5cm, align=center},
    arrow/.style={-{Stealth[length=3mm]}, thick}
]
    \node[level, fill=purple!20] (L4) at (0,6) {\textbf{Level 4: Civilizational} ($\Civilizational$)\\Historical trajectories, existential engineering};
    \node[level, fill=blue!20] (L3) at (0,4) {\textbf{Level 3: Cultural} ($\Cultural$)\\Overton windows, memetic warfare, hyperstition};
    \node[level, fill=green!20] (L2) at (0,2) {\textbf{Level 2: Social} ($\Social$)\\Group fractals, cascades, network dynamics};
    \node[level, fill=yellow!20] (L1) at (0,0) {\textbf{Level 1: Individual} ($\Individual$)\\Thought patterns, attractors, intervention};

    \draw[arrow] (L1.north) -- (L2.south) node[midway,right] {aggregates};
    \draw[arrow] (L2.north) -- (L3.south) node[midway,right] {aggregates};
    \draw[arrow] (L3.north) -- (L4.south) node[midway,right] {aggregates};

    \draw[arrow, dashed] (L4.west) to[out=180,in=180] node[midway,left] {constrains} (L3.west);
    \draw[arrow, dashed] (L3.west) to[out=180,in=180] node[midway,left] {constrains} (L2.west);
    \draw[arrow, dashed] (L2.west) to[out=180,in=180] node[midway,left] {constrains} (L1.west);
\end{tikzpicture}
\end{center}

\begin{definition}[The Cognitive Warfare Stack]
\label{def:cw-stack}
The \emph{TKS Cognitive Warfare Stack} is a graded category $\mathbf{CWS}$ with:
\begin{enumerate}[label=(\roman*)]
    \item \textbf{Objects}: Four levels $\Level{1}, \Level{2}, \Level{3}, \Level{4}$
    \item \textbf{Upward morphisms}: Aggregation functors $\Aggregate{i}{i+1} : \Level{i} \to \Level{i+1}$
    \item \textbf{Downward morphisms}: Constraint functors $\Project{i+1}{i} : \Level{i+1} \to \Level{i}$
    \item \textbf{Composition}: $\Aggregate{j}{k} \circ \Aggregate{i}{j} = \Aggregate{i}{k}$ for $i < j < k$
\end{enumerate}
subject to the \emph{cross-level coherence condition}:
\[
    \Project{i+1}{i} \circ \Aggregate{i}{i+1} \neq \mathrm{id}_{\Level{i}}
\]
(aggregation loses information; constraints don't recover it)
\end{definition}

\section{Level 1: Individual Psychology ($\Level{1} = \Individual$)}

Level 1 operates on the consciousness of individual persons, employing the
core TKS mathematical machinery developed in previous documents.

\begin{definition}[Level 1 Objects]
\label{def:L1-objects}
The objects of $\Level{1}$ are:
\begin{enumerate}[label=(\alph*)]
    \item \textbf{Psychological States}: Points $\psi \in \mathcal{S}_\psi$ in the
          consciousness state space
    \item \textbf{Attractors}: Stable configurations $\PsychAttractor \subset \mathcal{S}_\psi$
    \item \textbf{Fractals}: Thought-pattern iterators $\Frac : \mathcal{S}_\psi \to \mathcal{S}_\psi$
    \item \textbf{D/W/P Structures}: Acquisition chains $(D_m, W_m, P_m)$ for each
          Foundation $\Foundations_m$
\end{enumerate}
\end{definition}

\begin{definition}[Level 1 Capabilities]
\label{def:L1-capabilities}
The morphisms (operations) at Level 1 are:

\textbf{(a) Thought Pattern Extraction:}
\begin{equation}
\boxed{
    \mathsf{Extract} : \mathrm{Observable} \to \Frac
}
\end{equation}
Given observable behavior/communication, extract the governing fractal:
\[
    \mathsf{Extract}(\text{behavior}_1, \ldots, \text{behavior}_n) =
    \langle k_1 : k_2 : \cdots : k_m \rangle
\]

\textbf{(b) Attractor Prediction:}
\begin{equation}
\boxed{
    \mathsf{Predict} : (\psi_0, \Frac) \to \PsychAttractor
}
\end{equation}
Given initial state and governing fractal, predict the attractor:
\[
    \mathsf{Predict}(\psi_0, \Frac) = \lim_{n \to \infty} \Frac^n(\psi_0) = \PsychAttractor
\]

\textbf{(c) Precision Intervention:}
\begin{equation}
\boxed{
    \mathsf{Intervene} : (\PsychAttractor_{\mathrm{old}}, \PsychAttractor_{\mathrm{new}})
    \to \Frac_{\mathrm{intervention}}
}
\end{equation}
Compute the manipulation fractal that transforms one attractor to another:
\[
    \mathsf{Intervene}(X^*, Y^*) = \Frac_{\mathrm{destab}} \circ \Frac_{\mathrm{guide}}
    \circ \Frac_{\mathrm{stab}}
\]
\end{definition}

\begin{proposition}[Level 1 Core Equations]
\label{prop:L1-equations}
The fundamental equations governing Level 1 dynamics:

\textbf{Attractor Dynamics:}
\begin{equation}
    \frac{d\psi}{dt} = -\nabla V_{\mathrm{psych}}(\psi) + \eta(t)
\end{equation}
where $V_{\mathrm{psych}}$ is the psychological potential landscape with minima at attractors.

\textbf{Fractal Evolution:}
\begin{equation}
    \psi_{n+1} = \Frac(\psi_n) = \noe{k_m} \circ \cdots \circ \noe{k_1}(\psi_n)
\end{equation}

\textbf{Intervention Success Probability:}
\begin{equation}
    P(\PsychAttractor_{\mathrm{old}} \to \PsychAttractor_{\mathrm{new}}) =
    \exp\left(-\frac{\Delta V}{\sigma^2}\right) \cdot \prod_{A \in \mathrm{prereq}} \sigma(A)
\end{equation}
where $\Delta V$ is the potential barrier and $\sigma(A)$ are acquisition satisfaction indicators.
\end{proposition}

\section{Level 2: Social Dynamics ($\Level{2} = \Social$)}

Level 2 aggregates individual psychology into group phenomena, modeling how
individual fractals compose into collective dynamics.

\begin{definition}[Level 2 Objects]
\label{def:L2-objects}
The objects of $\Level{2}$ are:
\begin{enumerate}[label=(\alph*)]
    \item \textbf{Group States}: Distributions $\rho : \mathcal{S}_\psi \to [0,1]$ over
          individual states
    \item \textbf{Group Fractals}: Collective patterns $\GroupFrac$ governing group dynamics
    \item \textbf{Network Structures}: Graphs $G = (V, E)$ encoding social connections
    \item \textbf{Cascade States}: Propagation configurations $\Cascade \subset V \times [0,1]$
\end{enumerate}
\end{definition}

\begin{definition}[Level 2 Capabilities]
\label{def:L2-capabilities}

\textbf{(a) Group Fractal Modeling:}
\begin{equation}
\boxed{
    \GroupFrac = \mathbf{Agg}\left(\{\Frac_i\}_{i \in \mathrm{Group}}\right)
}
\end{equation}
Aggregate individual fractals into emergent group pattern:
\[
    \GroupFrac = \frac{1}{|G|} \sum_{i \in G} w_i \Frac_i +
    \sum_{(i,j) \in E} \lambda_{ij} (\Frac_i \otimes \Frac_j)
\]
where the cross-terms capture emergent effects from social coupling.

\textbf{(b) Cascade Optimization:}
\begin{equation}
\boxed{
    \mathsf{Cascade} : (\Memetic, G, \{v_{\mathrm{seed}}\}) \to \Cascade(t)
}
\end{equation}
Given a meme $\Memetic$, network $G$, and seed nodes, compute propagation:
\[
    \frac{d\Cascade_v}{dt} = \sum_{u \in N(v)} \alpha_{uv} \Cascade_u (1 - \Cascade_v)
    - \beta \Cascade_v
\]

\textbf{(c) Network Vulnerability Mapping:}
\begin{equation}
\boxed{
    \Vulnerability : G \to \{v_1^*, v_2^*, \ldots, v_k^*\}
}
\end{equation}
Identify optimal intervention points (influence maximizers):
\[
    \Vulnerability(G) = \arg\max_{S \subset V, |S| = k} \mathbb{E}[\Cascade_{\infty}(S)]
\]
\end{definition}

\begin{proposition}[Level 2 Core Equations]
\label{prop:L2-equations}

\textbf{Aggregation from Level 1:}
\begin{equation}
    \Aggregate{1}{2} : \{\psi_i, \Frac_i, \PsychAttractor_i\}_{i \in G} \mapsto
    (\rho_G, \GroupFrac, \mathcal{A}_G)
\end{equation}
where $\mathcal{A}_G$ is the collective attractor of the group.

\textbf{Group Attractor:}
\begin{equation}
    \mathcal{A}_G = \left\{ \rho : \GroupFrac(\rho) = \rho \right\}
\end{equation}

\textbf{Cascade Threshold:}
\begin{equation}
    \theta_{\mathrm{cascade}} = \frac{\beta}{\langle k \rangle \cdot \alpha_{\mathrm{avg}}}
\end{equation}
Cascade spreads globally when initial adoption exceeds $\theta_{\mathrm{cascade}}$.
\end{proposition}

\section{Level 3: Cultural Reality ($\Level{3} = \Cultural$)}

Level 3 aggregates social dynamics into cultural structures: the shared
reality-frames, acceptable discourse windows, and collective belief systems.

\begin{definition}[Level 3 Objects]
\label{def:L3-objects}
The objects of $\Level{3}$ are:
\begin{enumerate}[label=(\alph*)]
    \item \textbf{Overton Windows}: Intervals $\Overton \subset \mathcal{P}$ in
          policy/belief space defining ``acceptable'' discourse
    \item \textbf{Memetic Systems}: Reproducing idea-complexes $\Memetic$ with
          fitness functions
    \item \textbf{Hyperstitions}: Self-fulfilling belief structures $\Hyperstition$
    \item \textbf{Reality Frames}: Interpretive schemas $\mathcal{R}$ structuring perception
\end{enumerate}
\end{definition}

\begin{definition}[Level 3 Capabilities]
\label{def:L3-capabilities}

\textbf{(a) Overton Window Engineering:}
\begin{equation}
\boxed{
    \mathsf{Shift} : (\Overton_{\mathrm{old}}, \delta) \to \Overton_{\mathrm{new}}
}
\end{equation}
Shift the window of acceptable discourse:
\[
    \Overton_{\mathrm{new}} = \Overton_{\mathrm{old}} + \delta \cdot
    \mathrm{sign}\left(\sum_{\Memetic \in \mathrm{edge}} \mathrm{weight}(\Memetic)\right)
\]
Strategy: Introduce extreme positions to make moderate shifts seem reasonable.

\textbf{(b) Memetic Warfare Systems:}
\begin{equation}
\boxed{
    \frac{d\Memetic}{dt} = r_\Memetic \cdot \Memetic \cdot \left(1 - \frac{\Memetic}{K}\right)
    - \sum_{\Memetic' \in \mathrm{competing}} \alpha_{\Memetic\Memetic'} \Memetic \Memetic'
}
\end{equation}
Memes compete for mindshare following Lotka-Volterra dynamics.

\textbf{(c) Hyperstition Creation Engines:}
\begin{equation}
\boxed{
    \Hyperstition = (\mathcal{I}, \mathcal{B}, \mathcal{M}, \Frac_H)
    \quad \text{with } \frac{d\mathcal{B}}{dt} = f(\mathcal{M}(\mathcal{B}))
}
\end{equation}
Create self-fulfilling belief structures:
\[
    \frac{d\mathcal{B}}{dt} = \alpha \cdot \mathcal{M}(\mathcal{B}) - \beta \cdot \mathcal{B}
    + \gamma \cdot \mathrm{Propagation}(\mathcal{I})
\]
\end{definition}

\begin{proposition}[Level 3 Core Equations]
\label{prop:L3-equations}

\textbf{Aggregation from Level 2:}
\begin{equation}
    \Aggregate{2}{3} : \{\GroupFrac_G, \Cascade_G, \Vulnerability_G\}_{G \in \mathrm{Society}}
    \mapsto (\Overton, \{\Memetic_j\}, \{\Hyperstition_k\})
\end{equation}

\textbf{Overton Window Dynamics:}
\begin{equation}
    \frac{d\Overton}{dt} = \int_{\partial\Overton} \Memetic(x) \cdot n(x) \, dx
\end{equation}
where $\partial\Overton$ is the boundary and $n(x)$ is the outward normal.

\textbf{Reality Frame Competition:}
\begin{equation}
    P(\mathcal{R}_i \text{ adopted}) = \frac{e^{\beta U_i(\mathcal{R}_i)}}
    {\sum_j e^{\beta U_j(\mathcal{R}_j)}}
\end{equation}
where $U_i$ is the utility of frame $\mathcal{R}_i$ (explanatory power, emotional resonance, social proof).
\end{proposition}

\section{Level 4: Civilizational Scale ($\Level{4} = \Civilizational$)}

Level 4 operates on the longest time scales, modeling civilizational trajectories,
existential dynamics, and multi-generational attractor structures.

\begin{definition}[Level 4 Objects]
\label{def:L4-objects}
The objects of $\Level{4}$ are:
\begin{enumerate}[label=(\alph*)]
    \item \textbf{Historical Trajectories}: Paths $\Trajectory : [0,T] \to \mathcal{S}_{\mathrm{civ}}$
          through civilizational state space
    \item \textbf{Civilizational Attractors}: Long-term stable configurations
          $\mathcal{A}_{\mathrm{civ}}$
    \item \textbf{Existential Risks}: Probability measures $\mu_{\mathrm{risk}}$ over
          catastrophic events
    \item \textbf{Teleological Vectors}: Directional tendencies $\vec{\tau}$ in
          civilizational development
\end{enumerate}
\end{definition}

\begin{definition}[Level 4 Capabilities]
\label{def:L4-capabilities}

\textbf{(a) Historical Trajectory Modeling:}
\begin{equation}
\boxed{
    \Trajectory(t) = \int_0^t \vec{v}(\Trajectory(s), \Cultural(s)) \, ds
}
\end{equation}
Model civilizational evolution as flow through state space:
\[
    \frac{d\Trajectory}{dt} = F(\Trajectory) + G(\Trajectory, \Cultural) + \xi(t)
\]
where $F$ is endogenous dynamics, $G$ is cultural coupling, and $\xi$ is stochastic forcing.

\textbf{(b) Civilizational Attractor Computation:}
\begin{equation}
\boxed{
    \mathcal{A}_{\mathrm{civ}} = \lim_{t \to \infty} \Trajectory(t)
    \quad \text{(if limit exists)}
}
\end{equation}
Identify stable civilizational configurations:
\[
    \mathcal{A}_{\mathrm{civ}} = \{s \in \mathcal{S}_{\mathrm{civ}} :
    \nabla V_{\mathrm{civ}}(s) = 0, \; H(V_{\mathrm{civ}})(s) > 0\}
\]

\textbf{(c) Existential Risk/Reward Engineering:}
\begin{equation}
\boxed{
    \frac{d\mu_{\mathrm{risk}}}{dt} = -\nabla \cdot (v \mu_{\mathrm{risk}})
    + D \nabla^2 \mu_{\mathrm{risk}} + S
}
\end{equation}
Evolve probability distributions over existential outcomes:
\[
    P(\text{catastrophe in } [t, t+dt]) = \mu_{\mathrm{risk}}(\Trajectory(t)) \cdot dt
\]
\end{definition}

\begin{proposition}[Level 4 Core Equations]
\label{prop:L4-equations}

\textbf{Aggregation from Level 3:}
\begin{equation}
    \Aggregate{3}{4} : \{\Overton_j, \Memetic_j, \Hyperstition_j\}_{j \in \mathrm{Cultures}}
    \mapsto (\Trajectory, \mathcal{A}_{\mathrm{civ}}, \mu_{\mathrm{risk}})
\end{equation}

\textbf{Civilizational Potential:}
\begin{equation}
    V_{\mathrm{civ}}(\Trajectory) = -\int \left[ U(\Trajectory) - \lambda \cdot \mu_{\mathrm{risk}}(\Trajectory) \right] dt
\end{equation}
where $U$ is civilizational utility and $\lambda$ weights existential risk.

\textbf{Multi-Generational Fractal:}
\begin{equation}
    \TransFrac{\omega}(\mathrm{Civ}_0) = \lim_{\alpha \to \omega} \TransFrac{\alpha}(\mathrm{Civ}_0)
\end{equation}
The transfinite fractal limit captures civilizational convergence beyond finite time.
\end{proposition}

%% ============================================================================
%% CHAPTER: FUNCTORIAL RELATIONSHIPS
%% ============================================================================
\chapter{Functorial Relationships Between Levels}
\label{ch:functorial}

\section{The Stack as a Fibered Category}

\begin{definition}[Cognitive Stack Fibration]
\label{def:stack-fibration}
The Cognitive Warfare Stack forms a fibration $\pi : \mathbf{CWS} \to \mathbf{4}$ over
the ordinal category $\mathbf{4} = \{1 \to 2 \to 3 \to 4\}$:
\begin{itemize}
    \item The fiber $\pi^{-1}(i) = \Level{i}$ is the category of Level $i$ structures
    \item Vertical morphisms preserve level
    \item Cartesian lifts implement constraint propagation
\end{itemize}
\end{definition}

\section{Aggregation Functors}

\begin{theorem}[Aggregation Preserves Structure]
\label{thm:agg-structure}
Each aggregation functor $\Aggregate{i}{i+1}$ preserves:
\begin{enumerate}[label=(\roman*)]
    \item \textbf{Attractor structure}: If $\PsychAttractor$ is an attractor at Level $i$,
          then $\Aggregate{i}{i+1}(\PsychAttractor)$ is an attractor at Level $i+1$
    \item \textbf{Fractal composition}: $\Aggregate{i}{i+1}(\Frac_1 \circ \Frac_2) =
          \Aggregate{i}{i+1}(\Frac_1) \circ \Aggregate{i}{i+1}(\Frac_2)$
    \item \textbf{Dynamical flow}: If $\phi_t$ is a flow at Level $i$, then
          $\Aggregate{i}{i+1}(\phi_t)$ is a flow at Level $i+1$
\end{enumerate}
\end{theorem}

\begin{proof}[Proof Sketch]
(i) Follows from the limit-preservation property of aggregation: if individual states
converge to attractors, their distribution must converge to a distributional attractor.

(ii) Aggregation is a linear operation on fractals (weighted sum), and composition
distributes over addition.

(iii) The aggregated flow is the pushforward of the product flow by the aggregation map.
\end{proof}

\section{Constraint Functors}

\begin{theorem}[Constraints Impose Boundaries]
\label{thm:constraint-bounds}
Each constraint functor $\Project{i+1}{i}$ satisfies:
\begin{enumerate}[label=(\roman*)]
    \item \textbf{Boundary imposition}: $\Project{i+1}{i}(\Overton) \subset \mathcal{S}_\psi$
          defines the ``allowable'' psychological states
    \item \textbf{Information loss}: $\Project{i+1}{i} \circ \Aggregate{i}{i+1} \neq \mathrm{id}$
          (constraints don't recover individual information)
    \item \textbf{Attractor selection}: Higher-level attractors constrain which lower-level
          attractors are accessible
\end{enumerate}
\end{theorem}

\section{Scale Invariance (Self-Similarity)}

\begin{theorem}[Fractal Scale Invariance]
\label{thm:scale-invariance}
The same mathematical structures appear at all four levels:
\begin{equation}
    \Level{i} \cong \Level{j} \quad \text{(as categories)}
\end{equation}
with the isomorphism given by rescaling:
\begin{itemize}
    \item $\PsychAttractor_{\Level{1}} \mapsto \mathcal{A}_G$ (individual $\to$ group attractor)
    \item $\mathcal{A}_G \mapsto \Overton$ (group $\to$ cultural attractor)
    \item $\Overton \mapsto \mathcal{A}_{\mathrm{civ}}$ (cultural $\to$ civilizational attractor)
\end{itemize}
\end{theorem}

This scale invariance is the mathematical expression of the ``fractal'' nature of
consciousness organization: the same patterns recur at every scale, with the same
mathematical operations applicable.

%% ============================================================================
%% PLAIN ENGLISH EXPLANATION
%% ============================================================================
\chapter{Plain English Explanation}
\label{ch:plain-english}

\begin{tcolorbox}[colback=green!5!white,colframe=green!65!black,title=\textbf{The Control Tower Metaphor},breakable]

\textbf{Imagine a Control Tower with Four Altitude Levels}

You're in an air traffic control tower that can zoom between four different altitude
views. Each altitude shows the same airspace, but reveals different patterns:

\bigskip
\textbf{Ground Level (Level 1: Individual Psychology)}

At ground level, you see individual planes---their heading, speed, fuel, mechanical
status. You can:
\begin{itemize}
    \item \textbf{Extract patterns}: ``This pilot always banks left before climbing''
    \item \textbf{Predict destinations}: ``Given this flight path, they'll land at JFK''
    \item \textbf{Intervene precisely}: ``If we radio now, we can redirect them''
\end{itemize}

\textbf{TKS Translation}: Individual thought patterns are like flight paths. Psychological
attractors are like preferred destinations. Intervention fractals are like radio
instructions that change the course.

\bigskip
\textbf{Regional View (Level 2: Social Dynamics)}

Zooming up, you see traffic patterns---clusters of planes, congestion zones, cascade
effects when one plane delays others. You can:
\begin{itemize}
    \item \textbf{Model group behavior}: ``The 8am rush from Chicago creates this pattern''
    \item \textbf{Optimize cascades}: ``Delaying these 3 planes prevents 50 delays''
    \item \textbf{Map vulnerabilities}: ``This hub is the critical chokepoint''
\end{itemize}

\textbf{TKS Translation}: Group fractals emerge from individual patterns. Social cascades
spread like traffic congestion. Network vulnerabilities are the critical hubs.

\bigskip
\textbf{National View (Level 3: Cultural Reality)}

Higher still, you see airline policies, flight corridors, the ``acceptable'' routes that
everyone agrees on. You can:
\begin{itemize}
    \item \textbf{Engineer windows}: ``Opening this corridor changes what's 'normal'''
    \item \textbf{Run memetic warfare}: ``Our 'safety first' campaign beats their 'efficiency' message''
    \item \textbf{Create hyperstitions}: ``If everyone believes this route is best, it becomes best''
\end{itemize}

\textbf{TKS Translation}: Overton windows are like accepted flight corridors. Memes compete
like airline marketing. Hyperstitions are beliefs that become true because everyone acts
on them.

\bigskip
\textbf{Global View (Level 4: Civilizational Scale)}

At maximum altitude, you see decades of air traffic evolution---how routes developed,
where civilization is heading, existential risks to the entire system. You can:
\begin{itemize}
    \item \textbf{Model trajectories}: ``Aviation evolved from this point toward this attractor''
    \item \textbf{Compute attractors}: ``All paths lead to these stable configurations''
    \item \textbf{Engineer risks}: ``This intervention prevents catastrophic failure''
\end{itemize}

\textbf{TKS Translation}: Historical trajectories are civilizational flight paths.
Attractors are the stable endpoints. Existential risks are the crashes that end everything.

\bigskip
\textbf{Key Insight: Same Math, Different Scales}

The mathematics is \emph{identical} at all four levels. A ``fractal'' at Level 1
(individual thought pattern) has the same mathematical form as a ``fractal'' at Level 4
(civilizational trajectory). This is \textbf{scale invariance}---the hallmark of fractal systems.

What changes is:
\begin{itemize}
    \item \textbf{Timescale}: Seconds (L1) $\to$ Hours (L2) $\to$ Years (L3) $\to$ Centuries (L4)
    \item \textbf{Population}: 1 person $\to$ Groups $\to$ Cultures $\to$ Civilizations
    \item \textbf{Intervention difficulty}: Easy $\to$ Hard $\to$ Harder $\to$ Nearly impossible
\end{itemize}

But the \emph{equations} are the same. This is why understanding Level 1 deeply
teaches you Level 4---they're the same patterns at different scales.

\end{tcolorbox}

%% ============================================================================
%% PART II: PER-LEVEL SPECIFICATIONS
%% ============================================================================
\part{Per-Level Technical Specifications}

\chapter{Level 1: Individual Psychology}
\label{ch:level1-spec}

\section{Core TKS Mathematics Used}

\begin{center}
\renewcommand{\arraystretch}{1.4}
\begin{tabularx}{\textwidth}{>{\raggedright\arraybackslash}p{3cm}|>{\raggedright\arraybackslash}X|>{\raggedright\arraybackslash}p{4cm}}
\toprule
\textbf{TKS Component} & \textbf{Application} & \textbf{Key Equation} \\
\midrule
Noetic Operators $\noe{0}$--$\noe{9}$ & Classify mental operations & $\psi' = \noe{k}(\psi)$ \\
Noetic Fractals $\Frac$ & Model thought patterns & $\Frac = \langle k_1 : \cdots : k_n \rangle$ \\
Psychological Attractors & Predict stable states & $\PsychAttractor = \lim_{n\to\infty} \Frac^n(\psi_0)$ \\
D/W/P Chains & Map prerequisites & $D_m \to W_m \to P_m$ \\
RPM Monad & Compute with prerequisites & $\RPM(X) = (D \to W \to P) \times X$ \\
Manipulation Fractals & Design interventions & $\ManipFrac : \PsychAttractor_1 \to \PsychAttractor_2$ \\
\bottomrule
\end{tabularx}
\end{center}

\section{Input/Output Specification}

\begin{tcolorbox}[colback=blue!5!white,colframe=blue!65!black,title=\textbf{Level 1 I/O Specification}]
\textbf{INPUTS:}
\begin{itemize}
    \item Observable behavior sequences $\{b_1, b_2, \ldots, b_n\}$
    \item Communication/language samples
    \item Emotional state indicators
    \item Decision history
    \item Environmental context
\end{itemize}

\textbf{PROCESSING:}
\begin{enumerate}
    \item Map observables to World/Element coordinates: $b_i \mapsto W_n^k_{mf}$
    \item Extract governing fractal: $\{W_n^k\} \mapsto \Frac$
    \item Compute attractor: $\Frac \mapsto \PsychAttractor$
    \item Evaluate D/W/P structure: $\{(D_m, W_m, P_m)\}$
    \item Identify vulnerabilities: $V = \{m : \sigma(W_m) = \mathtt{false}\}$
\end{enumerate}

\textbf{OUTPUTS:}
\begin{itemize}
    \item Predicted attractor states with confidence intervals
    \item Optimal intervention fractals ranked by efficacy
    \item Vulnerability map identifying leverage points
    \item Timeline predictions with probability distributions
\end{itemize}
\end{tcolorbox}

\section{Interaction with Level 2}

\textbf{Upward (to Level 2):}
\begin{equation}
    \Aggregate{1}{2}(\{\psi_i, \Frac_i, \PsychAttractor_i\}) =
    \left(\rho = \frac{1}{N}\sum_i \delta_{\psi_i}, \;
    \GroupFrac = \sum_i w_i \Frac_i + \sum_{ij} \lambda_{ij} \Frac_i \otimes \Frac_j\right)
\end{equation}

\textbf{Downward (from Level 2):}
\begin{equation}
    \Project{2}{1}(\GroupFrac, G) = \{\text{peer influence on individual } i\} =
    \sum_{j \in N(i)} \alpha_{ij} \Frac_j
\end{equation}

%% ============================================================================
\chapter{Level 2: Social Dynamics}
\label{ch:level2-spec}

\section{Core TKS Mathematics Used}

\begin{center}
\renewcommand{\arraystretch}{1.4}
\begin{tabularx}{\textwidth}{>{\raggedright\arraybackslash}p{3cm}|>{\raggedright\arraybackslash}X|>{\raggedright\arraybackslash}p{4cm}}
\toprule
\textbf{TKS Component} & \textbf{Application} & \textbf{Key Equation} \\
\midrule
Group Fractals & Emergent collective patterns & $\GroupFrac = \sum w_i \Frac_i + \text{cross-terms}$ \\
Cascade Dynamics & Information/behavior spread & $\frac{d\Cascade}{dt} = \alpha(1-\Cascade) - \beta\Cascade$ \\
Network Graphs & Social structure & $G = (V, E, w)$ \\
Influence Centrality & Key node identification & $c_i = \sum_j A_{ij} c_j / \lambda_1$ \\
Threshold Models & Collective action & $\Cascade_i(t+1) = \mathbf{1}[\sum_j w_{ij}\Cascade_j > \theta_i]$ \\
\bottomrule
\end{tabularx}
\end{center}

\section{Input/Output Specification}

\begin{tcolorbox}[colback=blue!5!white,colframe=blue!65!black,title=\textbf{Level 2 I/O Specification}]
\textbf{INPUTS:}
\begin{itemize}
    \item Individual-level data from Level 1: $\{\psi_i, \Frac_i, \PsychAttractor_i\}$
    \item Network structure: adjacency matrix, edge weights
    \item Communication flows between nodes
    \item Group membership and identity markers
\end{itemize}

\textbf{PROCESSING:}
\begin{enumerate}
    \item Construct network graph: $G = (V, E, w)$
    \item Compute centrality measures: eigenvector, betweenness, closeness
    \item Aggregate individual fractals: $\GroupFrac$
    \item Simulate cascade dynamics: $\Cascade(t)$
    \item Identify optimal seed sets: $\arg\max_S \mathbb{E}[\Cascade_\infty(S)]$
\end{enumerate}

\textbf{OUTPUTS:}
\begin{itemize}
    \item Cascade reach predictions for different seed sets
    \item Vulnerability ranking of network nodes
    \item Group attractor predictions
    \item Optimal intervention timing windows
\end{itemize}
\end{tcolorbox}

\section{Interaction with Adjacent Levels}

\textbf{Upward (to Level 3):}
\begin{equation}
    \Aggregate{2}{3}(\{\GroupFrac_G\}_{G \in \mathrm{Society}}) =
    \Overton[\text{aggregate of group positions}]
\end{equation}

\textbf{Downward (from Level 3):}
\begin{equation}
    \Project{3}{2}(\Overton, \Memetic) = \text{allowable group positions} \subset \mathcal{P}
\end{equation}

%% ============================================================================
\chapter{Level 3: Cultural Reality}
\label{ch:level3-spec}

\section{Core TKS Mathematics Used}

\begin{center}
\renewcommand{\arraystretch}{1.4}
\begin{tabularx}{\textwidth}{>{\raggedright\arraybackslash}p{3cm}|>{\raggedright\arraybackslash}X|>{\raggedright\arraybackslash}p{4cm}}
\toprule
\textbf{TKS Component} & \textbf{Application} & \textbf{Key Equation} \\
\midrule
Overton Dynamics & Acceptable discourse bounds & $\frac{d\Overton}{dt} = \int_{\partial\Overton} \Memetic \cdot n \, dx$ \\
Memetic Ecology & Competing idea systems & $\frac{d\Memetic}{dt} = r\Memetic(1-\Memetic/K) - \alpha\Memetic\Memetic'$ \\
Hyperstition Loops & Self-fulfilling beliefs & $\frac{d\mathcal{B}}{dt} = \alpha\mathcal{M}(\mathcal{B}) - \beta\mathcal{B}$ \\
Reality Frames & Interpretive schemas & $P(\mathcal{R}_i) \propto e^{\beta U_i}$ \\
Transfinite Fractals & Long-term cultural evolution & $\TransFrac{\omega}(\Cultural_0)$ \\
\bottomrule
\end{tabularx}
\end{center}

\section{Input/Output Specification}

\begin{tcolorbox}[colback=blue!5!white,colframe=blue!65!black,title=\textbf{Level 3 I/O Specification}]
\textbf{INPUTS:}
\begin{itemize}
    \item Level 2 group data: $\{\GroupFrac_G, \Cascade_G\}$
    \item Media content and discourse samples
    \item Polling/sentiment data across populations
    \item Historical cultural trajectory data
\end{itemize}

\textbf{PROCESSING:}
\begin{enumerate}
    \item Estimate current Overton window: $\Overton(t)$
    \item Map memetic ecosystem: $\{\Memetic_j, r_j, K_j, \alpha_{jk}\}$
    \item Identify active hyperstitions: $\{\Hyperstition_k\}$
    \item Compute reality frame distribution: $P(\mathcal{R}_i)$
    \item Simulate cultural evolution: $\Cultural(t + \Delta t)$
\end{enumerate}

\textbf{OUTPUTS:}
\begin{itemize}
    \item Overton window shift predictions
    \item Optimal meme injection points and content
    \item Hyperstition viability assessments
    \item Cultural trajectory forecasts
\end{itemize}
\end{tcolorbox}

%% ============================================================================
\chapter{Level 4: Civilizational Scale}
\label{ch:level4-spec}

\section{Core TKS Mathematics Used}

\begin{center}
\renewcommand{\arraystretch}{1.4}
\begin{tabularx}{\textwidth}{>{\raggedright\arraybackslash}p{3cm}|>{\raggedright\arraybackslash}X|>{\raggedright\arraybackslash}p{4cm}}
\toprule
\textbf{TKS Component} & \textbf{Application} & \textbf{Key Equation} \\
\midrule
Trajectory Modeling & Civilizational paths & $\frac{d\Trajectory}{dt} = F(\Trajectory) + G(\Cultural)$ \\
Civilizational Attractors & Stable endpoints & $\mathcal{A}_{\mathrm{civ}} = \{s : \nabla V = 0\}$ \\
Existential Measures & Risk distributions & $\frac{d\mu}{dt} = -\nabla\cdot(v\mu) + D\nabla^2\mu$ \\
Teleological Vectors & Directional tendencies & $\vec{\tau} = \mathbb{E}[\Trajectory'(t)]$ \\
$\omega$-Fractals & Transfinite iteration & $\TransFrac{\omega} = \lim_{\alpha\to\omega}\TransFrac{\alpha}$ \\
\bottomrule
\end{tabularx}
\end{center}

\section{Input/Output Specification}

\begin{tcolorbox}[colback=blue!5!white,colframe=blue!65!black,title=\textbf{Level 4 I/O Specification}]
\textbf{INPUTS:}
\begin{itemize}
    \item Level 3 cultural data: $\{\Overton_j, \Memetic_j, \Hyperstition_j\}$
    \item Historical time series (centuries of data)
    \item Cross-civilizational comparative data
    \item Technological capability trajectories
\end{itemize}

\textbf{PROCESSING:}
\begin{enumerate}
    \item Fit trajectory model: $\Trajectory(t)$ parameters
    \item Compute civilizational potential landscape: $V_{\mathrm{civ}}$
    \item Identify attractor basins: $\{\mathcal{A}_k, \mathrm{Basin}_k\}$
    \item Estimate risk probabilities: $\mu_{\mathrm{risk}}(\Trajectory)$
    \item Compute teleological vectors: $\vec{\tau}$
\end{enumerate}

\textbf{OUTPUTS:}
\begin{itemize}
    \item Long-term trajectory predictions with uncertainty bands
    \item Attractor identification and basin mapping
    \item Existential risk forecasts
    \item Optimal civilizational intervention strategies
\end{itemize}
\end{tcolorbox}

%% ============================================================================
%% PART III: CAPABILITY MATRIX
%% ============================================================================
\part{Capability Matrix and Technology Mapping}

\chapter{Capability Matrix}
\label{ch:capability-matrix}

\section{Core Capabilities Table}

\begin{center}
\small
\renewcommand{\arraystretch}{1.5}
\begin{tabularx}{\textwidth}{>{\raggedright\arraybackslash}p{2.5cm}|>{\raggedright\arraybackslash}p{3cm}|c|c|c}
\toprule
\textbf{Capability} & \textbf{TKS Technology} & \textbf{Accuracy} & \textbf{Scale} & \textbf{Stealth} \\
\midrule
Thought Prediction & Fractal extraction + attractor computation & High (0.7--0.9) & L1 & Medium \\
\hline
Intent Detection & D/W/P chain analysis & High (0.8--0.95) & L1 & High \\
\hline
Emotional State & C-World mapping + oscillation detection & Medium (0.6--0.8) & L1 & High \\
\hline
Group Behavior & Group fractal + cascade modeling & Medium (0.5--0.7) & L2 & High \\
\hline
Action Timeline & RPM prerequisite evaluation & Medium (0.5--0.7) & L1--L2 & Medium \\
\hline
Vulnerability ID & Lacunarity + missing prerequisite detection & High (0.75--0.9) & L1--L3 & High \\
\hline
Manipulation Efficacy & Intervention fractal + barrier analysis & Variable (0.4--0.8) & L1--L4 & Low--High \\
\bottomrule
\end{tabularx}
\end{center}

\section{Capability Specifications}

\subsection{Thought Prediction}

\begin{tcolorbox}[colback=yellow!5!white,colframe=yellow!65!black,title=\textbf{Thought Prediction Capability}]
\textbf{TKS Technology Stack:}
\begin{enumerate}
    \item \textbf{Fractal Extraction}: Observe behavior sequence, map to noetic operators,
          identify dominant fractal $\Frac = \langle k_1 : \cdots : k_n \rangle$
    \item \textbf{Attractor Computation}: Iterate fractal to find $\PsychAttractor = \lim \Frac^n(\psi_0)$
    \item \textbf{Prediction Generation}: Future thoughts lie in basin of attraction
\end{enumerate}

\textbf{Accuracy Factors:}
\begin{itemize}
    \item Higher with more observation data
    \item Higher for stable attractors (low lacunarity)
    \item Lower for strange attractors (chaotic patterns)
    \item Degraded by conscious deception or pattern-breaking
\end{itemize}

\textbf{Stealth Level}: Medium --- requires observation but not direct interaction
\end{tcolorbox}

\subsection{Intent Detection}

\begin{tcolorbox}[colback=yellow!5!white,colframe=yellow!65!black,title=\textbf{Intent Detection Capability}]
\textbf{TKS Technology Stack:}
\begin{enumerate}
    \item \textbf{D-Chain Analysis}: Identify activated Desires $D_1, \ldots, D_7$
    \item \textbf{W-Chain Status}: Check Wisdom acquisition for each desire
    \item \textbf{P-Chain Prediction}: Infer likely Power manifestations
\end{enumerate}

\textbf{Key Equation:}
\begin{equation}
    P(\text{Intent}_m) = \sigma(D_m) \times P(\sigma(W_m) | D_m) \times \mathsf{Context}
\end{equation}

\textbf{Accuracy Factors:}
\begin{itemize}
    \item Highest accuracy for strong desires (high $D_m$ activation)
    \item Can detect intent before action (D precedes P)
    \item May misattribute Foundation (desire for power vs. wealth)
\end{itemize}

\textbf{Stealth Level}: High --- infers from observable behavior without intrusion
\end{tcolorbox}

\subsection{Vulnerability Identification}

\begin{tcolorbox}[colback=yellow!5!white,colframe=yellow!65!black,title=\textbf{Vulnerability Identification Capability}]
\textbf{TKS Technology Stack:}
\begin{enumerate}
    \item \textbf{Lacunarity Mapping}: Identify high-lacunarity regions (gaps)
    \item \textbf{Prerequisite Audit}: Find missing acquisitions $\{A : \sigma(A) = \mathtt{false}\}$
    \item \textbf{Attractor Basin Analysis}: Identify unstable regions in psychological landscape
\end{enumerate}

\textbf{Vulnerability Types:}
\begin{center}
\small
\begin{tabular}{l|l|l}
\textbf{Type} & \textbf{Detection Method} & \textbf{Intervention Vector} \\
\hline
Cognitive Gap & $\noe{1}$ deficit & Information injection \\
Emotional Instability & Limit cycle attractor & Resonance amplification \\
Identity Fragility & High $\Lacunarity(\Frac_{\mathrm{self}})$ & Identity challenge \\
Desire-Wisdom Gap & $D_m$ high, $W_m$ low & Offer false wisdom \\
Social Isolation & Low network centrality & Offer belonging \\
\end{tabular}
\end{center}

\textbf{Stealth Level}: High --- purely analytical, no target interaction
\end{tcolorbox}

%% ============================================================================
%% PART IV: ETHICAL BOUNDARIES
%% ============================================================================
\part{Ethical Boundaries Framework}

\chapter{Ethical Boundaries Framework}
\label{ch:ethics}

\begin{tcolorbox}[colback=gray!10!white,colframe=gray!75!black,title=\textbf{Framing Note}]
The capabilities documented in this manual exist independently of ethical frameworks.
They have been used throughout history by states, corporations, religions, and
individuals. The question is not ``should these tools exist?'' but ``given that
they exist, what boundaries should govern their use?''

This chapter documents the ethical terrain, not to advocate for any position,
but to make the considerations explicit.
\end{tcolorbox}

\section{Red Lines: Non-Negotiable Prohibitions}

\begin{tcolorbox}[colback=red!10!white,colframe=red!75!black,title=\textbf{Category I: Red Lines},breakable]

These represent violations of fundamental autonomy that cannot be ethically justified
under any framework that values persons as ends rather than means:

\bigskip
\textbf{1. Non-Consensual Mind Alteration}
\begin{equation}
\boxed{
    \ManipFrac : \PsychAttractor_{\mathrm{other}} \to \PsychAttractor'
    \quad \text{WITHOUT } \mathrm{Consent} = \mathtt{PROHIBITED}
}
\end{equation}
Applying manipulation fractals to alter another person's psychological attractors
without their informed consent constitutes a violation equivalent to physical assault
on consciousness.

\textbf{Boundary Condition:}
\[
    \forall \ManipFrac_{\mathrm{other}} : \mathrm{Consent}(\mathrm{target}, \ManipFrac) \in
    \{\mathtt{explicit}, \mathtt{implied\_emergency}\}
\]

\bigskip
\textbf{2. Reality Tunneling}
\begin{equation}
\boxed{
    \Hyperstition_{\mathrm{coercive}} : \mathcal{R}_{\mathrm{target}} \to \mathcal{R}_{\mathrm{controlled}}
    \quad \text{= PROHIBITED}
}
\end{equation}
Deliberately constructing hyperstitions or reality frames that trap targets in
controlled perception systems (cults, closed ideological systems, manufactured
dependencies).

\textbf{Boundary Condition:}
\[
    \forall \Hyperstition : \exists \text{ exit pathway } \land \text{target knows it}
\]

\bigskip
\textbf{3. Agency Erosion}
\begin{equation}
\boxed{
    \mathrm{Intervention} : \mathrm{Agency}_{\mathrm{target}} \to \mathrm{Agency}' < \mathrm{Agency}
    \quad \text{= PROHIBITED}
}
\end{equation}
Any intervention that systematically reduces the target's capacity for autonomous
choice, self-reflection, or independent action.

\textbf{Boundary Condition:}
\[
    \frac{d(\mathrm{Agency}_{\mathrm{target}})}{d(\mathrm{Intervention})} \geq 0
\]
Interventions must preserve or enhance agency, never diminish it.

\end{tcolorbox}

\section{Gray Areas: Context-Dependent Ethics}

\begin{tcolorbox}[colback=gray!20!white,colframe=gray!60!black,title=\textbf{Category II: Gray Areas},breakable]

These represent domains where ethical assessment depends heavily on context,
intent, and outcome. The same technique may be ethical or unethical depending
on circumstances:

\bigskip
\textbf{1. Therapeutic vs. Manipulative}

\begin{center}
\begin{tabular}{p{6cm}|p{6cm}}
\textbf{Therapeutic (Generally Ethical)} & \textbf{Manipulative (Generally Unethical)} \\
\hline
Client seeks change & Target unaware of influence \\
Explicit goals agreed upon & Goals hidden or misrepresented \\
Increases client's self-understanding & Decreases target's self-understanding \\
Client can terminate at will & Target cannot easily exit \\
Practitioner has fiduciary duty & Actor has competing interests \\
\end{tabular}
\end{center}

\textbf{The Boundary Equation:}
\begin{equation}
    \mathrm{Ethics}(\ManipFrac) = f(\mathrm{Consent}, \mathrm{Transparency},
    \mathrm{Benefit}_{\mathrm{target}}, \mathrm{Exit\_Option})
\end{equation}

\bigskip
\textbf{2. Education vs. Indoctrination}

\begin{center}
\begin{tabular}{p{6cm}|p{6cm}}
\textbf{Education (Generally Ethical)} & \textbf{Indoctrination (Generally Unethical)} \\
\hline
Teaches how to think & Teaches what to think \\
Encourages questioning & Punishes questioning \\
Exposes to multiple perspectives & Restricts to single perspective \\
Develops critical faculties & Suppresses critical faculties \\
Student can reject teachings & Rejection has severe consequences \\
\end{tabular}
\end{center}

\textbf{The Boundary Equation:}
\begin{equation}
    \mathrm{Ethics}(\Cultural_{\mathrm{transmission}}) =
    f(\mathrm{Pluralism}, \mathrm{Critical\_Space}, \mathrm{Exit\_Cost})
\end{equation}

\bigskip
\textbf{3. Persuasion vs. Coercion}

The spectrum from legitimate persuasion to illegitimate coercion:

\begin{center}
\begin{tikzpicture}
    \draw[thick,->] (0,0) -- (12,0);
    \node at (0,-0.5) {Coercion};
    \node at (12,-0.5) {Persuasion};

    \draw (2,0.2) -- (2,-0.2) node[below] {\small Threats};
    \draw (4,0.2) -- (4,-0.2) node[below] {\small Deception};
    \draw (6,0.2) -- (6,-0.2) node[below] {\small Manipulation};
    \draw (8,0.2) -- (8,-0.2) node[below] {\small Influence};
    \draw (10,0.2) -- (10,-0.2) node[below] {\small Information};

    \fill[red!30] (0,0.5) rectangle (5,1);
    \fill[yellow!30] (5,0.5) rectangle (8,1);
    \fill[green!30] (8,0.5) rectangle (12,1);

    \node at (2.5,0.75) {\tiny Unethical};
    \node at (6.5,0.75) {\tiny Context-dep.};
    \node at (10,0.75) {\tiny Ethical};
\end{tikzpicture}
\end{center}

\end{tcolorbox}

\section{Historical Precedents}

\begin{tcolorbox}[colback=blue!5!white,colframe=blue!65!black,title=\textbf{Pre-TKS Analogues},breakable]

The capabilities documented in this manual are formalizations of techniques that
have existed throughout history. TKS provides mathematical language; it does not
create new capabilities:

\bigskip
\textbf{Level 1 (Individual) Precedents:}
\begin{itemize}
    \item Rhetoric (Aristotle, 4th c. BCE) --- systematic persuasion techniques
    \item Confession/Spiritual Direction (Catholic tradition) --- psychological mapping
    \item Psychoanalysis (Freud, 1890s) --- attractor identification and intervention
    \item Advertising psychology (20th c.) --- desire-based manipulation
    \item Cognitive behavioral therapy --- fractal restructuring
\end{itemize}

\textbf{Level 2 (Social) Precedents:}
\begin{itemize}
    \item Political organizing (all eras) --- cascade optimization
    \item Propaganda ministries (20th c.) --- systematic social engineering
    \item Network marketing (1940s--) --- cascade design
    \item Social media algorithms (2010s--) --- computational cascade optimization
\end{itemize}

\textbf{Level 3 (Cultural) Precedents:}
\begin{itemize}
    \item Religious institutions (all eras) --- reality frame installation
    \item State education systems (19th c.--) --- cultural reproduction
    \item Public relations industry (1920s--) --- Overton engineering
    \item Think tanks and foundations (20th c.--) --- systematic memetic warfare
\end{itemize}

\textbf{Level 4 (Civilizational) Precedents:}
\begin{itemize}
    \item Imperial grand strategy (all eras) --- trajectory engineering
    \item Geopolitical theory (Mackinder, Spykman) --- civilizational modeling
    \item Long-term institutional planning (Vatican, etc.) --- multi-generational attractors
\end{itemize}

\textbf{Key Insight}: TKS does not create these capabilities; it makes them
mathematically explicit. The ethical questions are ancient; only the precision
is new.

\end{tcolorbox}

%% ============================================================================
%% PART V: APPLICATION SCENARIOS
%% ============================================================================
\part{Full-Stack Application Scenarios}

\chapter{Full-Stack Applications}
\label{ch:applications}

\begin{tcolorbox}[colback=yellow!5!white,colframe=yellow!65!black,title=\textbf{Scenario Framework}]
Each scenario demonstrates how all four levels of the stack coordinate to produce
integrated effects. These are presented as technical descriptions without ethical
endorsement---documenting what is possible, not what should be done.
\end{tcolorbox}

\section{Scenario 1: Full-Stack Therapeutic Intervention}

\textbf{Context}: Individual with treatment-resistant depression, social isolation,
internalized stigma, in a culture that pathologizes mental health treatment.

\begin{tcolorbox}[colback=green!5!white,colframe=green!65!black,title=\textbf{Therapeutic Stack Deployment},breakable]

\textbf{Level 1 --- Individual Psychology:}
\begin{itemize}
    \item \textbf{Assessment}: Extract governing fractal $\Frac_{\mathrm{dep}} = \langle 3:5:3:5 \rangle$
          (rejection-absorption-rejection-absorption loop)
    \item \textbf{Attractor ID}: $\PsychAttractor = C3^5_{3c}$ (absorbed rejection in emotional world)
    \item \textbf{Vulnerability}: High lacunarity in $\Foundations_4$ (Companionship) ---
          $\sigma(W_4) = \mathtt{false}$
    \item \textbf{Intervention}: Apply $\ManipFrac_{\mathrm{therapy}} = \langle 8:1:4:2:5:9 \rangle$
          (cause $\to$ awareness $\to$ energy $\to$ positive $\to$ receive $\to$ complete)
\end{itemize}

\textbf{Level 2 --- Social Dynamics:}
\begin{itemize}
    \item \textbf{Network Analysis}: Identify isolation pattern, map remaining connections
    \item \textbf{Support Cascade}: Coordinate family/friends to provide consistent support
    \item \textbf{Group Integration}: Connect to peer support group (cascade seed for belonging)
    \item \textbf{Equation}: $\Cascade_{\mathrm{support}}(v_{\mathrm{target}}) = \sum_{u \in \mathrm{support}} \alpha_{uv}$
\end{itemize}

\textbf{Level 3 --- Cultural Reality:}
\begin{itemize}
    \item \textbf{Stigma Assessment}: Map cultural frame $\mathcal{R}_{\mathrm{stigma}}$
    \item \textbf{Counter-Narrative}: Introduce alternative frame $\mathcal{R}_{\mathrm{recovery}}$
    \item \textbf{Validation Pathway}: Connect to cultural narratives of healing, growth, resilience
    \item \textbf{Equation}: $P(\mathcal{R}_{\mathrm{recovery}}) \uparrow$ via exposure to counter-examples
\end{itemize}

\textbf{Level 4 --- Civilizational Context:}
\begin{itemize}
    \item \textbf{Long-term Trajectory}: Understand depression epidemic in civilizational context
    \item \textbf{Systemic Factors}: Acknowledge social/economic drivers beyond individual
    \item \textbf{Meaning Frame}: Connect individual recovery to larger purpose
\end{itemize}

\textbf{Integration Equation:}
\begin{equation}
    \mathrm{Recovery} = \ManipFrac_{\mathrm{L1}} \times \Cascade_{\mathrm{L2}} \times
    \mathcal{R}_{\mathrm{L3}} \times \Trajectory_{\mathrm{L4}}
\end{equation}

\textbf{Outcome}: Individual attractor shifts from depression to euthymia, supported by
social network, validated by cultural frame, meaningful in civilizational context.

\end{tcolorbox}

\section{Scenario 2: Full-Stack Public Health Campaign}

\textbf{Context}: Vaccination campaign in population with vaccine hesitancy,
misinformation networks, and cultural distrust of institutions.

\begin{tcolorbox}[colback=green!5!white,colframe=green!65!black,title=\textbf{Public Health Stack Deployment},breakable]

\textbf{Level 1 --- Individual Psychology:}
\begin{itemize}
    \item \textbf{Hesitancy Mapping}: Identify dominant hesitancy attractors
    \item \textbf{Fractal Types}: Fear-based $\langle 3:4:3 \rangle$, distrust-based $\langle 3:1:3 \rangle$,
          identity-based $\langle 5:3:5 \rangle$
    \item \textbf{Intervention Design}: Tailored messaging per fractal type
    \item \textbf{D/W/P Analysis}: Address desire for safety ($D_3$), provide wisdom ($W_3$)
\end{itemize}

\textbf{Level 2 --- Social Dynamics:}
\begin{itemize}
    \item \textbf{Network Mapping}: Identify opinion leaders, community structure
    \item \textbf{Cascade Design}: Optimize seed selection for maximum reach
    \item \textbf{Counter-Cascade}: Identify and address misinformation super-spreaders
    \item \textbf{Social Proof}: Leverage peer effects for normalization
\end{itemize}

\textbf{Level 3 --- Cultural Reality:}
\begin{itemize}
    \item \textbf{Overton Analysis}: Map current window of vaccine acceptability
    \item \textbf{Frame Competition}: Counter misinformation frames with evidence-based frames
    \item \textbf{Hyperstition Design}: ``Vaccination is what responsible community members do''
    \item \textbf{Cultural Validation}: Connect to existing values (family protection, community care)
\end{itemize}

\textbf{Level 4 --- Civilizational Context:}
\begin{itemize}
    \item \textbf{Historical Trajectory}: Position within history of disease eradication
    \item \textbf{Collective Action Frame}: ``Previous generations eliminated smallpox; we can do this''
    \item \textbf{Long-term Stakes}: Connect to civilizational resilience against pandemics
\end{itemize}

\textbf{Integration:}
\begin{equation}
    \mathrm{Uptake} = \int \left[ \ManipFrac_{\mathrm{L1}}(x) \times \Cascade_{\mathrm{L2}}(x) \times
    \Overton_{\mathrm{L3}}(x) \right] d\mu_{\mathrm{population}}(x)
\end{equation}

\end{tcolorbox}

\section{Scenario 3: Full-Stack Counter-Terrorism}

\textbf{Context}: Preventing radicalization pipeline from online exposure to violent action.

\begin{tcolorbox}[colback=green!5!white,colframe=green!65!black,title=\textbf{Counter-Terrorism Stack Deployment},breakable]

\textbf{Level 1 --- Individual Psychology:}
\begin{itemize}
    \item \textbf{Vulnerability Detection}: High lacunarity in identity ($\Foundations_1$),
          belonging ($\Foundations_4$), power ($\Foundations_5$)
    \item \textbf{Radicalization Fractal}: $\langle 3:2:4:6 \rangle$
          (rejection of outgroup $\to$ attraction to ingroup $\to$ energization $\to$ action)
    \item \textbf{Intervention Window}: Before $\noe{6}$ (action/projection) crystallizes
    \item \textbf{Counter-Fractal}: $\langle 1:2:5:4 \rangle$
          (awareness $\to$ positive reframe $\to$ receive alternatives $\to$ redirect energy)
\end{itemize}

\textbf{Level 2 --- Social Dynamics:}
\begin{itemize}
    \item \textbf{Network Detection}: Map online/offline radicalization networks
    \item \textbf{Cascade Disruption}: Identify and address key propagation nodes
    \item \textbf{Alternative Networks}: Connect at-risk individuals to prosocial groups
    \item \textbf{Former Exit Support}: Leverage former radicals as counter-narrative sources
\end{itemize}

\textbf{Level 3 --- Cultural Reality:}
\begin{itemize}
    \item \textbf{Narrative Warfare}: Counter extremist hyperstitions with alternative meanings
    \item \textbf{Overton Management}: Maintain firm boundary against violence-legitimizing discourse
    \item \textbf{Identity Provision}: Offer compelling non-violent identity narratives
    \item \textbf{Grievance Address}: Where legitimate grievances exist, provide non-violent channels
\end{itemize}

\textbf{Level 4 --- Civilizational Context:}
\begin{itemize}
    \item \textbf{Root Cause Analysis}: Understand civilizational dynamics producing radicalization
    \item \textbf{Systemic Intervention}: Address underlying drivers (alienation, inequality, meaning-crisis)
    \item \textbf{Long-term Trajectory}: Build civilizational resilience against extremism
\end{itemize}

\textbf{Integration:}
\begin{equation}
    P(\text{prevention}) = \prod_{i=1}^{4} \left(1 - P(\text{radicalization} | \text{no } L_i \text{ intervention})\right)
\end{equation}

Each level provides defense-in-depth; failure at one level can be caught by another.

\end{tcolorbox}

%% ============================================================================
%% CONCLUSION
%% ============================================================================
\chapter*{Conclusion: The Stack Exists}
\addcontentsline{toc}{chapter}{Conclusion: The Stack Exists}

\begin{tcolorbox}[colback=gray!10!white,colframe=gray!75!black,title=\textbf{Final Technical Note}]

The TKS Cognitive Warfare Stack is not a proposal; it is a description. The
mathematical structures documented here formalize capabilities that have always
existed and have always been used. Advertising firms model psychological attractors.
Political campaigns optimize cascades. Think tanks engineer Overton windows.
Geopoliticians model civilizational trajectories.

TKS provides precise mathematical language for these operations. This precision
enables:
\begin{enumerate}
    \item \textbf{Analysis}: Understanding how influence actually works
    \item \textbf{Defense}: Recognizing when these techniques are being used
    \item \textbf{Optimization}: More effective deployment when deployment is justified
    \item \textbf{Accountability}: Clear standards for ethical evaluation
\end{enumerate}

The question is not whether to allow such capabilities to exist---they already do.
The question is whether understanding should be concentrated among elites or
distributed broadly. This document contributes to distribution.

\textbf{The mathematics is neutral. The ethics are not. Use wisely.}

\end{tcolorbox}

%% ============================================================================
%% APPENDIX: QUICK REFERENCE
%% ============================================================================
\appendix
\chapter{Quick Reference Tables}
\label{app:quick-reference}

\section{Level Summary}

\begin{center}
\renewcommand{\arraystretch}{1.5}
\begin{tabularx}{\textwidth}{c|X|X|c}
\toprule
\textbf{Level} & \textbf{Domain} & \textbf{Key Capabilities} & \textbf{Timescale} \\
\midrule
L1 & Individual Psychology & Thought prediction, attractor analysis, precision intervention & Seconds--Days \\
L2 & Social Dynamics & Group modeling, cascade optimization, network vulnerability & Hours--Months \\
L3 & Cultural Reality & Overton engineering, memetic warfare, hyperstition & Months--Decades \\
L4 & Civilizational & Trajectory modeling, attractor computation, existential engineering & Decades--Centuries \\
\bottomrule
\end{tabularx}
\end{center}

\section{Core Equations Summary}

\begin{center}
\renewcommand{\arraystretch}{1.8}
\begin{tabularx}{\textwidth}{l|X}
\toprule
\textbf{Equation} & \textbf{Meaning} \\
\midrule
$\PsychAttractor = \lim_{n\to\infty} \Frac^n(\psi_0)$ & Psychological attractor as fractal limit \\
$\GroupFrac = \sum_i w_i \Frac_i + \sum_{ij} \lambda_{ij} \Frac_i \otimes \Frac_j$ & Group fractal with emergent cross-terms \\
$\frac{d\Overton}{dt} = \int_{\partial\Overton} \Memetic \cdot n \, dx$ & Overton window dynamics \\
$\frac{d\Trajectory}{dt} = F(\Trajectory) + G(\Cultural) + \xi$ & Civilizational trajectory equation \\
$\Aggregate{i}{i+1} \circ \Project{i+1}{i} \neq \mathrm{id}$ & Information loss in level transitions \\
\bottomrule
\end{tabularx}
\end{center}

\section{Ethical Decision Matrix}

\begin{center}
\renewcommand{\arraystretch}{1.4}
\begin{tabularx}{\textwidth}{X|c|c|c}
\toprule
\textbf{Factor} & \textbf{Red Line} & \textbf{Gray Area} & \textbf{Generally OK} \\
\midrule
Consent & None & Implicit/Situational & Explicit \\
Transparency & Deceptive & Partial & Full \\
Target Benefit & Harm & Ambiguous & Clear Benefit \\
Exit Option & Trapped & Difficult & Easy \\
Agency Effect & Diminished & Unchanged & Enhanced \\
\bottomrule
\end{tabularx}
\end{center}

\end{document}
